% ============================================================
%  CHUONG 1. TONG QUAN
% ============================================================
\chapter{Tổng quan}
\label{ch:tong-quan}

\section{Bối cảnh và nhu cầu thực tiễn}
\label{sec:boi-canh}

Nấu ăn tại nhà là hoạt động thường nhật, nhưng khác với nhiều hoạt động thường
nhật khác, nó đòi hỏi một chuỗi quyết định có ràng buộc lẫn nhau. Quyết định nấu
món gì phụ thuộc vào nguyên liệu đang có. Nguyên liệu cần mua lại phụ thuộc vào
các món dự định nấu. Và thực đơn của một ngày nên xét trong tương quan với cả
tuần để tránh trùng lặp và bảo đảm đa dạng.

Xét theo góc độ tin học, chuỗi quyết định này có thể quy về ba bài toán có cấu
trúc rõ ràng:

\textbf{Bài toán 1: Truy vấn ngược từ nguyên liệu.} Cho trước tập nguyên liệu
$A$ mà người dùng đang có và một kho công thức, cần tìm và xếp hạng những món có
thể nấu được. Điểm đặc biệt của bài toán này là chiều truy vấn ngược với cách tổ
chức dữ liệu thông thường: các trang công thức thường được tổ chức để tra cứu
\textit{từ tên món ra nguyên liệu}, trong khi nhu cầu thực tế của người nội trợ
là tra cứu \textit{từ nguyên liệu ra tên món}.

\textbf{Bài toán 2: Xếp lịch có ràng buộc.} Cần gán món ăn vào 21 ô của một lưới
7 ngày $\times$ 3 bữa, thỏa mãn đồng thời nhiều ràng buộc: hạn chế lặp món, phù
hợp với từng bữa trong ngày, tôn trọng lựa chọn về vùng miền và cân đối năng
lượng theo ngày. Đây là một biến thể của bài toán xếp lịch (scheduling) có ràng
buộc.

\textbf{Bài toán 3: Gộp nhóm và tổng hợp.} Từ thực đơn đã lập, cần tổng hợp
nguyên liệu cần mua. Cùng một nguyên liệu có thể xuất hiện ở nhiều món, và cùng
một món có thể xuất hiện nhiều lần trong tuần, do đó phải gộp nhóm theo nguyên
liệu và cộng dồn khối lượng.

Cả ba bài toán đều thao tác trên dữ liệu có cấu trúc và đều có lời giải thuật
toán xác định, không đòi hỏi kỹ thuật trí tuệ nhân tạo phức tạp. Đây là cơ sở để
khẳng định tính khả thi của đề tài.

\section{Khảo sát các ứng dụng tương tự}
\label{sec:khao-sat}

Để xác định khoảng trống mà đồ án hướng tới, tác giả đã khảo sát các ứng dụng
nấu ăn và lập thực đơn tiêu biểu, cả quốc tế lẫn trong nước. Toàn bộ thông tin
dưới đây được tra cứu trực tiếp từ trang chủ hoặc trang phân phối chính thức của
từng sản phẩm trong ngày 25/07/2026. Nguyên tắc khảo sát được đặt ra ngay từ đầu
là chỉ ghi nhận những chức năng có thể kiểm chứng được từ nguồn chính chủ; với
những chức năng chỉ được nhắc tới trên nguồn thứ cấp, báo cáo ghi rõ là không xác
minh được thay vì suy đoán.

\subsection{Cookpad}

Cookpad là nền tảng chia sẻ công thức nấu ăn theo mô hình cộng đồng, trong đó nội
dung do chính người dùng đăng tải thay vì do một đội ngũ biên tập tạo ra [1].
Cookpad có phiên bản tiếng Việt riêng tại địa chỉ cookpad.com/vn, nằm trong hệ
thống hơn ba mươi phiên bản theo quốc gia [5].

Về chức năng, mô tả chính thức của ứng dụng cho biết Cookpad hỗ trợ \textbf{tìm
công thức theo nguyên liệu} với thông điệp ``nấu những bữa ăn ngon từ những gì
bạn đã có sẵn trong tủ lạnh'', đồng thời có khả năng \textbf{tạo danh sách đi chợ
tự động từ công thức} và cho phép người dùng \textbf{xây dựng thực đơn theo tuần}
[1].

\textit{Nhận xét.} Cookpad là sản phẩm gần với đề tài nhất về mặt chức năng. Tuy
nhiên mô tả chính thức không nêu rõ việc xây dựng thực đơn tuần là \textit{sinh
tự động} hay chỉ là công cụ để người dùng \textit{tự sắp xếp} công thức vào lịch,
nên không thể khẳng định Cookpad có thuật toán sinh thực đơn tự động. Hạn chế
đáng kể hơn nằm ở mô hình nội dung: vì công thức do người dùng tự đăng nên chất
lượng và độ chính xác về định lượng nguyên liệu không đồng đều, điều này ảnh
hưởng trực tiếp tới độ tin cậy của danh sách đi chợ được sinh ra.

\subsection{Yummly}

Yummly từng là một dịch vụ gợi ý công thức được nhắc tới nhiều trong các nghiên
cứu về hệ khuyến nghị ẩm thực. Tuy nhiên, kết quả kiểm tra trực tiếp trong quá
trình khảo sát cho thấy \textbf{dịch vụ này đã ngừng hoạt động}: truy cập
\url{https://www.yummly.com/} hiện bị chuyển hướng vĩnh viễn (HTTP 301) sang
trang công thức của KitchenAid, và tên miền help.yummly.com không còn phân giải
được [15]. Theo nguồn thứ cấp trích lại thông báo của công ty chủ quản, website
và ứng dụng di động của Yummly đã đóng cửa từ ngày 20/12/2024 [14].

\textit{Nhận xét.} Trường hợp Yummly có giá trị tham chiếu quan trọng đối với đề
tài theo hai hướng. Thứ nhất, nó cho thấy người dùng Việt Nam không thể trông cậy
vào tính sẵn sàng lâu dài của các dịch vụ gợi ý ẩm thực quốc tế, vốn phụ thuộc
vào quyết định kinh doanh của doanh nghiệp nước ngoài. Thứ hai, do dịch vụ đã
đóng, mọi mô tả chi tiết về tính năng của Yummly đều không thể kiểm chứng lại ở
thời điểm hiện tại; báo cáo này vì vậy không đưa ra bất kỳ khẳng định nào về các
chức năng cụ thể của sản phẩm.

\subsection{Mealime}

Mealime là ứng dụng chuyên biệt cho việc \textbf{lập thực đơn}, không phải nền
tảng chia sẻ công thức. Trang chủ của sản phẩm nêu khả năng lập kế hoạch bữa ăn
cho cả tuần dựa trên tùy chọn cá nhân về khẩu phần, chế độ ăn và dị ứng thực
phẩm, đồng thời \textbf{tự động tổng hợp nguyên liệu thành danh sách đi chợ theo
danh mục} [9]. Sản phẩm áp dụng mô hình freemium với gói Mealime Pro giá 2,99 USD
mỗi tháng [2].

\textit{Nhận xét.} Mealime giải quyết tốt bài toán lập thực đơn và danh sách đi
chợ, nhưng chức năng tìm món \textbf{xuất phát từ nguyên liệu người dùng đang có}
không được nêu trong mô tả chính thức; trọng tâm của sản phẩm là sinh thực đơn từ
sở thích và chế độ ăn. Ngoài ra, Mealime không có nội dung tiếng Việt và không có
món ăn Việt Nam, nên khó áp dụng trực tiếp cho bữa cơm gia đình Việt.

\subsection{Paprika Recipe Manager}

Paprika là ứng dụng quản lý công thức cá nhân đa nền tảng, cho phép lưu công thức
từ web, lập kế hoạch bữa ăn và tạo danh sách đi chợ với khả năng gộp nguyên liệu
trùng nhau và sắp xếp theo khu vực gian hàng [13]. Ứng dụng hỗ trợ tìm kiếm theo
tên và theo nguyên liệu [3].

\textit{Nhận xét.} Điểm khác biệt căn bản so với đề tài là Paprika \textbf{không
sinh thực đơn tự động}: người dùng phải tự kéo thả từng công thức vào lịch [13].
Nói cách khác, Paprika số hóa thao tác lập thực đơn thủ công chứ không thay người
dùng đưa ra quyết định.

\subsection{Các website nấu ăn trong nước}

Ở trong nước, Esheep Kitchen là một blog công thức nấu ăn tiếng Việt có lượng
người theo dõi lớn, với nội dung xoay quanh món Việt, bánh và đồ uống [7]. Kiểm
tra trực tiếp cho thấy đây thuần túy là nền tảng đăng tải nội dung: trang không
có chức năng tìm món theo nguyên liệu sẵn có, không sinh thực đơn tuần và không
xuất danh sách đi chợ [7].

Một trường hợp đáng chú ý khác là Cooky.vn. Trang này khởi đầu năm 2015 như một
website công thức nấu ăn, sau đó từ năm 2020 chuyển hướng sang mô hình đi chợ hộ
và giao thực phẩm tươi. Đến tháng 11/2023, doanh nghiệp đã thu hẹp hoạt động, rút
khỏi thị trường Hà Nội và chỉ còn vận hành tại Thành phố Hồ Chí Minh [4]. Tại
thời điểm khảo sát, tác giả \textbf{không truy cập được} trang cooky.vn do chứng
chỉ bảo mật TLS của trang đã hết hạn, vì vậy báo cáo này không đưa ra kết luận nào
về các chức năng hiện có của sản phẩm.

\subsection{Bảng tổng hợp kết quả khảo sát}

Bảng~\ref{tab:khao-sat} tổng hợp kết quả khảo sát. Các ô ghi ``Không xác minh
được'' phản ánh việc tác giả không tìm được nguồn chính thức để khẳng định, chứ
không đồng nghĩa với việc chức năng đó chắc chắn không tồn tại.

\begin{table}[H]
\centering
\caption{So sánh chức năng giữa các ứng dụng khảo sát và đề tài}
\label{tab:khao-sat}
\begin{tabularx}{\textwidth}{p{1.9cm} Y Y Y p{1.5cm} Y}
\toprule
\bfseries Ứng dụng & \bfseries Gợi ý theo nguyên liệu sẵn có &
\bfseries Sinh thực đơn tuần tự động & \bfseries Xuất danh sách đi chợ &
\bfseries Nội dung tiếng Việt & \bfseries Mô hình giá \\
\midrule
Cookpad & Có [1] & Có, chưa rõ mức độ tự động [1] & Có [1] & Có [5] &
  Miễn phí kèm gói Premium [1] \\
\rowhl Yummly & Đã ngừng hoạt động từ 20/12/2024 [14][15] & Đã ngừng hoạt động &
  Đã ngừng hoạt động & Đã ngừng hoạt động & Đã ngừng hoạt động \\
Mealime & Không nêu trong mô tả chính thức [9][2] & Có [9] & Có [9] &
  Không xác minh được & Freemium, Pro 2,99 USD/tháng [2] \\
\rowhl Paprika & Có [3] & Không, lập kế hoạch thủ công [13] & Có [13] &
  Không xác minh được & Trả phí một lần [13] \\
Esheep Kitchen & Không [7] & Không [7] & Không [7] & Có [7] & Miễn phí [7] \\
\rowhl Cooky.vn & Không xác minh được & Không xác minh được &
  Không xác minh được & Có [4] & Không xác minh được \\
\midrule
\bfseries Cooking\-Advisor & \bfseries Có & \bfseries Có & \bfseries Có &
  \bfseries Có & \bfseries Miễn phí \\
\bottomrule
\end{tabularx}
\end{table}

\section{Nhận xét chung và khoảng trống nghiên cứu}
\label{sec:khoang-trong}

Từ kết quả khảo sát, có thể rút ra bốn nhận xét.

\textbf{Một là, ba chức năng cốt lõi hiếm khi xuất hiện đồng thời trong một sản
phẩm.} Các sản phẩm khảo sát có xu hướng mạnh ở một khâu và yếu ở khâu còn lại:
Mealime mạnh về lập thực đơn nhưng không hỗ trợ truy vấn ngược từ nguyên liệu;
Paprika mạnh về quản lý công thức và danh sách đi chợ nhưng không tự động hóa
việc lập thực đơn; các blog trong nước chỉ dừng ở việc cung cấp nội dung.

\textbf{Hai là, các sản phẩm quốc tế không phù hợp với bối cảnh ẩm thực Việt
Nam.} Mealime và Paprika không có nội dung tiếng Việt và không có dữ liệu món ăn
Việt. Ngay cả khi giao diện được dịch, cơ sở dữ liệu nguyên liệu và công thức vẫn
xây dựng theo khẩu vị Âu - Mỹ, khiến kết quả gợi ý không sát với thực tế bữa cơm
gia đình Việt.

\textbf{Ba là, tính sẵn sàng lâu dài của dịch vụ nước ngoài là một rủi ro thực
tế.} Việc Yummly ngừng hoạt động hoàn toàn vào cuối năm 2024 là minh chứng cụ thể
cho rủi ro này [14][15].

\textbf{Bốn là, mô hình nội dung do cộng đồng đóng góp ảnh hưởng tới chất lượng
dữ liệu định lượng.} Với các nền tảng như Cookpad, công thức do người dùng tự
đăng nên định lượng nguyên liệu không được chuẩn hóa. Điều này không gây trở ngại
khi người dùng chỉ đọc công thức, nhưng trở thành vấn đề khi hệ thống cần
\textit{tính toán} trên dữ liệu đó, ví dụ khi cộng dồn khối lượng để sinh danh
sách đi chợ.

Khoảng trống mà đồ án hướng tới, vì vậy, là một hệ thống \textbf{tích hợp cả ba
chức năng trong một luồng công việc liền mạch}, xây dựng trên \textbf{dữ liệu món
ăn Việt Nam được chuẩn hóa về định lượng}, với thuật toán xử lý minh bạch và có
thể kiểm chứng được.

\section{Vấn đề tập trung giải quyết và đóng góp của đồ án}
\label{sec:van-de}

Trên cơ sở phân tích trên, đồ án tập trung giải quyết bài toán sau:

\begin{quote}
\textit{Cho một kho dữ liệu món ăn Việt Nam được chuẩn hóa, hãy xây dựng một hệ
thống web cho phép:}

\textit{(a) xếp hạng các món có thể nấu từ một tập nguyên liệu cho trước;}

\textit{(b) sinh tự động thực đơn cho bảy ngày với các ràng buộc về tính đa dạng
và cân đối năng lượng;}

\textit{(c) tổng hợp danh sách nguyên liệu cần mua từ thực đơn đó.}
\end{quote}

Đóng góp của đồ án gồm bốn điểm:

\begin{enumerate}[label=\arabic*.]
  \item \textbf{Về mô hình dữ liệu:} thiết kế lược đồ quan hệ chuẩn hóa cho bài
        toán, trong đó bảng trung gian giữa món ăn và nguyên liệu mang thêm
        thuộc tính định lượng và đơn vị. Đây là điều kiện tiên quyết để danh sách
        đi chợ có thể cộng dồn khối lượng một cách chính xác.
  \item \textbf{Về thuật toán gợi ý:} đề xuất sử dụng \textbf{độ phủ} thay cho
        hệ số Jaccard thông dụng, kèm lập luận về lý do lựa chọn (trình bày ở mục
        \ref{sec:do-do}), cùng cơ chế xếp hạng nhiều tiêu chí giúp kết quả ổn
        định.
  \item \textbf{Về thuật toán lập thực đơn:} xây dựng thuật toán tham lam có ràng
        buộc theo thứ tự nới lỏng xác định, trong đó ràng buộc vùng miền được giữ
        cứng còn ràng buộc về bữa ăn chỉ được nới sau khi đã cạn nguồn món chưa
        dùng.
  \item \textbf{Về kiểm chứng:} xây dựng bộ kiểm thử đơn vị tự động cho cả ba
        thuật toán, qua đó phát hiện được một khiếm khuyết thực tế về tính đa
        dạng của thực đơn mà quan sát bằng mắt thường khó nhận ra.
\end{enumerate}
